% ============================================================
\section{11.1 RNN: 은닉 상태와 순전파}
\begin{frame}{왜 순서가 중요한가}
  ``강아지가 고양이를 쫓는다''와 ``고양이가 강아지를 쫓는다''는
  \textbf{똑같은 단어 세 개}로 이루어져 있지만 완전히 다른 문장이다.
  \begin{itemize}
    \item MLP·CNN에 이 문장을 넣으려면 단어를 \textbf{벡터 하나로
      뭉쳐야} 한다 $\to$ \textbf{순서 정보가 사라짐}
    \item 주가 시계열: 같은 100개 숫자라도 순서에 따라 ``상승 후
      조정''이든 ``조정 후 상승''이든 다른 신호
    \item DNA·단백질: 같은 염기 네 개(A, C, G, T)의 \textbf{배열}
      자체가 의미를 결정
    \item CNN(Week 10)은 공간적 인접성 + 위치 불변성을 활용
    \item 시퀀스의 ``왼쪽$\to$오른쪽''은 \textbf{시간의 방향성}:
      앞 단어는 뒤 단어를 영향주지만 뒤 단어는 앞 단어를 알 수 없다
  \end{itemize}
  \vskip0.5em
  \kq{이 비대칭을 모델이 반영할 수 있어야 한다 --- RNN의 출발 질문.}
\end{frame}
